\begin{exercises}
   \recommended \item \label{exer:IndNilLftShift}
      \definend{右移}算子的幂零指数是多少？这里它作用在实数三元组构成的空间上：
       \begin{equation*}
          (x,y,z)\mapsto(0,x,y)
       \end{equation*}
       \begin{answer}
         是三。
         说它至少是三，是因为 $\ell^2(\,(1,1,1)\,)=(0,0,1)\neq \zero$。
         说它至多是三，是因为
         $(x,y,z)\mapsto (0,x,y)\mapsto (0,0,x)\mapsto (0,0,0)$。
       \end{answer}
   \recommended \item
     对下面每个链基，给出幂零映射的幂零指数，并给出该幂零映射每次迭代的
     值域空间与零空间的维数。
     \begin{exparts}
       \partsitem $
         \begin{strings}{ccccccc}
            \vec{\beta}_1 &\mapsto &\vec{\beta}_2 &\mapsto &\zero  \\
            \vec{\beta}_3 &\mapsto &\vec{\beta}_4 &\mapsto &\zero
          \end{strings}$
       \partsitem $
         \begin{strings}{ccccccc}
            \vec{\beta}_1 &\mapsto &\vec{\beta}_2 &\mapsto &\vec{\beta}_3
                 &\mapsto &\zero  \\
            \vec{\beta}_4 &\mapsto &\zero \\
            \vec{\beta}_5 &\mapsto &\zero \\
            \vec{\beta}_6 &\mapsto &\zero
          \end{strings}$
       \partsitem $
         \begin{strings}{ccccccccc}
            \vec{\beta}_1 &\mapsto &\vec{\beta}_2 &\mapsto &\vec{\beta}_3
                 &\mapsto &\zero
          \end{strings}$
     \end{exparts}
     并给出该矩阵的标准形。
     \begin{answer}
       \begin{exparts}
         \partsitem 定义域的维数是四。
           该映射的作用是：空间中任意向量
           $c_1\cdot \vec{\beta}_1+c_2\cdot \vec{\beta}_2
             +c_3\cdot \vec{\beta}_3+c_4\cdot \vec{\beta}_4$
           被送到
           $c_1\cdot \vec{\beta}_2+c_2\cdot \zero
             +c_3\cdot \vec{\beta}_4+c_4\cdot \zero
            =c_1\cdot \vec{\beta}_3+c_3\cdot\vec{\beta}_4$。
           映射的第一次作用把两个基向量 $\vec{\beta}_2$ 与
           $\vec{\beta}_4$ 送到零，
           因此零空间的维数为~二，值域空间的维数也为~二。
           第二次作用后，四个基向量全部被送到零，
           于是二次幂的零空间维数为~四，而二次幂的值域空间维数为零。
           因此幂零指数是二。
           这就是标准形。
           \begin{equation*}
             \begin{mat}[r]
               0  &0  &0  &0  \\
               1  &0  &0  &0  \\
               0  &0  &0  &0  \\
               0  &0  &1  &0
             \end{mat}
           \end{equation*}
         \partsitem 该映射定义域的维数是六。
           对一次幂而言，零空间的维数是四，值域空间的维数是二。
           对二次幂而言，零空间的维数是五，值域空间的维数是一。
           第三次迭代得到的零空间维数为六，值域空间维数为零。
           幂零指数是三，这就是标准形。
           \begin{equation*}
             \begin{mat}[r]
               0  &0  &0  &0  &0  &0  \\
               1  &0  &0  &0  &0  &0  \\
               0  &1  &0  &0  &0  &0  \\
               0  &0  &0  &0  &0  &0  \\
               0  &0  &0  &0  &0  &0  \\
               0  &0  &0  &0  &0  &0
             \end{mat}
           \end{equation*}
         \partsitem 定义域的维数是三，幂零指数是三。
           一次幂的零空间维数为一，值域空间维数为二。
           二次幂的零空间维数为二，值域空间维数为一。
           最后，三次幂的零空间维数为三，值域空间维数为零。
           这是标准形矩阵。
           \begin{equation*}
             \begin{mat}[r]
               0  &0  &0  \\
               1  &0  &0  \\
               0  &1  &0
             \end{mat}
          \end{equation*}
      \end{exparts}
    \end{answer}
  \item
    判断下列矩阵中哪些是幂零的。
    \begin{exparts*}
      \partsitem
        $\begin{mat}[r]
           -2  &4  \\
           -1  &2
         \end{mat}$
       \partsitem
         $\begin{mat}[r]
           3  &1  \\
           1  &3
         \end{mat}$
       \partsitem
         $\begin{mat}[r]
           -3  &2  &1  \\
           -3  &2  &1  \\
           -3  &2  &1
         \end{mat}$
       \partsitem
         $\begin{mat}[r]
            1  &1  &4  \\
            3  &0  &-1 \\
            5  &2  &7
         \end{mat}$
       \partsitem
         $\begin{mat}[r]
            45  &-22  &-19  \\
            33  &-16  &-14  \\
            69  &-34  &-29
         \end{mat}$
     \end{exparts*}
     \begin{answer}
       由 \nearbylemma{le:RangeAndNullChains}，零度会在第 $n$ 次迭代时增长到尽可能大，其中 $n$ 是定义域的维数。
       于是，对 $\nbyn{2}$ 矩阵，
       我们只需检查其平方是否为零矩阵。
       对 $\nbyn{3}$ 矩阵，只需检查其立方。
       \begin{exparts}
         \partsitem 是，该矩阵是幂零的，因为
           它的平方是零矩阵。
         \partsitem 否，它的平方不是零矩阵。
           \begin{equation*}
             \begin{mat}[r]
                3  &1  \\
                1  &3
             \end{mat}^2
             =\begin{mat}[r]
                10  &6  \\
                6   &10
             \end{mat}
           \end{equation*}
         \partsitem 是，它的立方是零矩阵。
           事实上，它的平方就是零。
         \partsitem 否，它的三次幂不是零矩阵。
           \begin{equation*}
             \begin{mat}[r]
               1  &1  &4  \\
               3  &0  &-1 \\
               5  &2  &7
             \end{mat}^3
             =\begin{mat}[r]
               206  &86  &304  \\
                26  &8   &26   \\
               438  &180 &634
             \end{mat}
           \end{equation*}
         \partsitem 是，该矩阵的立方是零矩阵。
       \end{exparts}
       要看出第二与第四个矩阵不是幂零的，另一种方法是注意它们是非奇异的。
     \end{answer}
   \recommended \item
     求该矩阵的标准形。
     \begin{equation*}
       \begin{mat}[r]
         0  &1  &1  &0  &1  \\
         0  &0  &1  &1  &1  \\
         0  &0  &0  &0  &0  \\
         0  &0  &0  &0  &0  \\
         0  &0  &0  &0  &0
       \end{mat}
     \end{equation*}
     \begin{answer} 计算表格
       \begin{center}
           \begin{tabular}{c|cc}
              \multicolumn{1}{c}{\( p \)}  &\( N^p \)  &\( \nullspace{N^p}  \) \\
              \hline
              \( 1 \)
                &\matrixvenlarge{\begin{mat}[r]
                    0  &1  &1  &0  &1  \\
                    0  &0  &1  &1  &1  \\
                    0  &0  &0  &0  &0  \\
                    0  &0  &0  &0  &0  \\
                    0  &0  &0  &0  &0
                   \end{mat}}
                &\( \set{\matrixvenlarge{\colvec{r \\ u \\ -u-v \\ u \\ v}}
                               \suchthat r,u,v\in\C}  \) \\
              \( 2 \)
                &\matrixvenlarge{\begin{mat}[r]
                    0  &0  &1  &1  &1  \\
                    0  &0  &0  &0  &0  \\
                    0  &0  &0  &0  &0  \\
                    0  &0  &0  &0  &0  \\
                    0  &0  &0  &0  &0
                   \end{mat}}
                &\( \set{\matrixvenlarge{\colvec{r \\ s \\ -u-v \\ u \\ v}}
                               \suchthat r,s,u,v\in\C}  \) \\
             \( 2 \)
                &\textit{--零矩阵--}
                &\( \C^5 \)
           \end{tabular}
         \end{center}
         由此得到对链基的如下要求：~有三个基向量直接被映射到零，
         另有一个基向量经第二次作用被映射到零，
         最后一个基向量
         则经第三次作用被映射到零。
         因此，链基具有如下形式。
         \begin{equation*}
           \begin{strings}{ccccccc}
             \vec{\beta}_1 &\mapsto &\vec{\beta}_2 &\mapsto
               &\vec{\beta}_3 &\mapsto &\zero               \\
             \vec{\beta}_4 &\mapsto &\zero                  \\
             \vec{\beta}_5 &\mapsto &\zero
           \end{strings}
         \end{equation*}
         据此立得标准形。
         \begin{equation*}
           \begin{mat}[r]
             0  &0  &0  &0  &0  \\
             1  &0  &0  &0  &0  \\
             0  &1  &0  &0  &0  \\
             0  &0  &0  &0  &0  \\
             0  &0  &0  &0  &0
           \end{mat}
         \end{equation*}
     \end{answer}
   \recommended \item
     考虑 \nearbyexample{ex:FiveByNilStrBas} 中的矩阵。
     \begin{exparts}
       \partsitem 利用该映射在链基上的作用给出标准形。
       \partsitem 求出把该矩阵化为标准形的换基矩阵。
       \partsitem 利用前一小题的答案检验第一小题的答案。
     \end{exparts}
     \begin{answer}
       \begin{exparts}
         \partsitem 标准形含一个 $\nbyn{3}$ 块与一个
           $\nbyn{2}$ 块
           \begin{equation*}
             \begin{pmat}{rrr|rr}
               0  &0  &0  &0  &0  \\
               1  &0  &0  &0  &0  \\
               0  &1  &0  &0  &0  \\ \hline
               0  &0  &0  &0  &0  \\
               0  &0  &0  &1  &0  \\
             \end{pmat}
           \end{equation*}
           分别对应于基中长度为三的链与长度为二的
           链。
         \partsitem 设 $N$ 是底层映射关于标准基的表示。
           设 $B$ 是我们即将换到的基。
           由相似图
           \begin{equation*}
             \begin{CD}
               \C^2_{\wrt{\stdbasis_2}}
                   @>n>N>
                   \C^2_{\wrt{\stdbasis_2}}     \\
              @V\scriptstyle\identity V\scriptstyle PV
                  @V\scriptstyle\identity V\scriptstyle PV \\
              C^2_{\wrt{B}}
                  @>n>>
              \C^2_{\wrt{B}}
            \end{CD}
          \end{equation*}
          可知标准形矩阵为 $PNP^{-1}$，其中
          \begin{equation*}
            P^{-1}
            =\rep{\identity}{B,\stdbasis_5}
            =\begin{mat}[r]
               1 &0 &0 &0 &0 \\
               0 &1 &0 &1 &0 \\
               1 &0 &1 &0 &0 \\
               0 &0 &1 &1 &1 \\
               0 &0 &0 &0 &1
             \end{mat}
          \end{equation*}
          而 $P$ 是它的逆。
          \begin{equation*}
            P=\rep{\identity}{\stdbasis_5,B}
             =(P^{-1})^{-1}
            =\begin{mat}[r]
               1 &0 &0 &0 &0 \\
              -1 &1 &1 &-1&1 \\
              -1 &0 &1 &0 &0 \\
               1 &0 &-1&1 &-1\\
               0 &0 &0 &0 &1
             \end{mat}
          \end{equation*}
         \partsitem 检验它的计算是例行的。
       \end{exparts}
     \end{answer}
   \recommended \item
     下列每个矩阵都是幂零的。
     \begin{exparts*}
       \partsitem \(
         \begin{mat}[r]
           1/2  &-1/2  \\
           1/2  &-1/2
         \end{mat}        \)
       \partsitem \(
         \begin{mat}[r]
           0  &0  &0  \\
           0  &-1 &1  \\
           0  &-1 &1
         \end{mat}        \)
       \partsitem \(
         \begin{mat}[r]
          -1  &1  &-1 \\
           1  &0  &1  \\
           1  &-1 &1
         \end{mat}        \)
     \end{exparts*}
     把每个矩阵化为标准形。
     \begin{answer}
       \begin{exparts}
       \partsitem 计算
         \begin{center}
           \begin{tabular}{c|cc}
              \multicolumn{1}{c}{\( p \)}  &\( N^p \)  &\( \nullspace{N^p} \) \\
             \hline
              \( 1 \)
                &\matrixvenlarge{\begin{mat}[r]
                     1/2  &-1/2  \\
                     1/2  &-1/2
                   \end{mat}}
                &\( \set{\matrixvenlarge{\colvec{u \\ u}}
                               \suchthat u\in\C}  \) \\
             \( 2 \)
                &\textit{--零矩阵--}
                &\( \C^2 \)
           \end{tabular}
         \end{center}
         表明由该矩阵表示的任何映射
         必须这样作用在链基上
         \begin{equation*}
           \begin{strings}{ccccccc}
             \vec{\beta}_1 &\mapsto &\vec{\beta}_2  &\mapsto &\zero
           \end{strings}
         \end{equation*}
         因为作用一次后零空间的维数为~一，
         且恰好有一个基向量 $\vec{\beta}_2$ 被映射到零。
         因此，这个关于
         $\sequence{\vec{\beta}_1,\vec{\beta}_2}$ 的表示就是标准形。
         \begin{equation*}
           \begin{mat}[r]
             0    &0   \\
             1    &0
           \end{mat}
         \end{equation*}
       \partsitem 这里的计算与前一题类似。
         \begin{center}
           \begin{tabular}{c|cc}
              \multicolumn{1}{c}{\( p \)}  &\( N^p \)  &\( \nullspace{N^p} \) \\
              \hline
              \( 1 \)
                &\matrixvenlarge{\begin{mat}[r]
                     0  &0  &0  \\
                     0  &-1 &1  \\
                     0  &-1 &1
                   \end{mat}}
                &\( \set{\matrixvenlarge{\colvec{u \\ v \\ v}}
                               \suchthat u,v\in\C}  \) \\
            \( 2 \)
                &\textit{--零矩阵--}
                &\( \C^3 \)
           \end{tabular}
        \end{center}
        表格表明链基形如
         \begin{equation*}
           \begin{strings}{ccccccc}
             \vec{\beta}_1 &\mapsto &\vec{\beta}_2 &\mapsto &\zero  \\
             \vec{\beta}_3 &\mapsto &\zero
           \end{strings}
         \end{equation*}
         因为映射作用一次后零空间的
         维数为二\Dash
         --$\vec{\beta}_2$ 与 $\vec{\beta}_3$ 都被送到零\Dash 而
         再迭代一次又使另一个向量
         被送到零。
       \partsitem 计算
         \begin{center}
           \begin{tabular}{c|cc}
              \multicolumn{1}{c}{\( p \)}  &\( N^p \)  &\( \nullspace{N^p} \) \\
              \hline
              \( 1 \)
                &\matrixvenlarge{\begin{mat}[r]
                     -1  &1  &-1  \\
                      1  &0  &1   \\
                      1  &-1 &1
                   \end{mat}}
                &\( \set{\matrixvenlarge{\colvec{u \\ 0 \\ -u}}
                               \suchthat u\in\C}  \) \\
              \( 2 \)
                &\matrixvenlarge{\begin{mat}[r]
                      1  &0  &1   \\
                      0  &0  &0   \\
                     -1  &0  &-1
                   \end{mat}}
                &\( \set{\matrixvenlarge{\colvec{u \\ v \\ -u}}
                               \suchthat u,v\in\C}  \) \\
             \( 3 \)
                &\textit{--零矩阵--}
                &\( \C^3 \)
           \end{tabular}
         \end{center}
         表明由该基表示的任何映射必须这样作用在
         链基上。
         \begin{equation*}
           \begin{strings}{ccccccc}
             \vec{\beta}_1 &\mapsto &\vec{\beta}_2 &\mapsto
               &\vec{\beta}_3 &\mapsto &\zero
           \end{strings}
         \end{equation*}
         因此，这就是标准形。
         \begin{equation*}
           \begin{mat}[r]
             0    &0   &0   \\
             1    &0   &0   \\
             0    &1   &0
           \end{mat}
         \end{equation*}
       \end{exparts}
      \end{answer}
   \item
     描述用幂零矩阵的标准形矩阵左乘或右乘所产生的效果。
     \begin{answer}
       下面几个例子
       \begin{equation*}
         \begin{mat}[r]
           0  &0  \\
           1  &0
         \end{mat}
         \begin{mat}
           a  &b  \\
           c  &d
         \end{mat}
         =
         \begin{mat}
           0  &0  \\
           a  &b
         \end{mat}
         \qquad
         \begin{mat}[r]
           0  &0  &0 \\
           1  &0  &0 \\
           0  &1  &0
         \end{mat}
         \begin{mat}
           a  &b  &c \\
           d  &e  &f \\
           g  &h  &i
         \end{mat}
         =
         \begin{mat}
           0  &0  &0 \\
           a  &b  &c \\
           d  &e  &f
         \end{mat}
       \end{equation*}
       表明用一块次对角线上的 $1$ 左乘一个矩阵，
       会把该矩阵的行向下移。
       而不同的块
       \begin{equation*}
         \begin{mat}[r]
           0  &0  &0  &0  \\
           1  &0  &0  &0  \\
           0  &0  &0  &0  \\
           0  &0  &1  &0
         \end{mat}
         \begin{mat}
           a  &b  &c  &d  \\
           e  &f  &g  &h  \\
           i  &j  &k  &l  \\
           m  &n  &o  &p
         \end{mat}
         =
         \begin{mat}
           0  &0  &0  &0  \\
           a  &b  &c  &d  \\
           0  &0  &0  &0  \\
           i  &j  &k  &l
         \end{mat}
       \end{equation*}
       则会把矩阵的不同部分向下移。

       右乘对列做的是类似的事情。
       参见 \nearbyexercise{exer:IndNilLftShift}。
     \end{answer}
   \item
     幂零性在相似下不变吗？
     也就是说，与幂零矩阵相似的矩阵必定也是幂零的吗？
     若是，指数也相同吗？
     \begin{answer}
       是。
       把 \nearbyexample{ex:NilMatNotCanon} 的最后一句推广即可。
       至于指数，同样的最后一句表明新矩阵的指数小于或等于 $\hat{N}$ 的指数，而把
       两个矩阵的角色互换便得到另一方向的不等式。

       这个问题的另一种回答方式是证明：一个矩阵是幂零的当且仅当与之相关的映射是幂零的，且
       指数相同。
       于是，由于相似矩阵表示同一个映射，结论
       随之而来。
       这就是下面的 \nearbyexercise{exer:MatNilIffMapNil}。
     \end{answer}
   \recommended \item
     证明幂零矩阵唯一的特征值是零。
     \begin{answer}
       注意标准形的幂零矩阵只有
       零特征值；例如，这个下三角矩阵
       \begin{equation*}
          \begin{mat}
            -x  &0  &0  \\
             1  &-x &0  \\
             0  &1  &-x
          \end{mat}
       \end{equation*}
       的行列式是 \( (-x)^3 \)，它的唯一的根是零。
       而相似矩阵有相同的特征值，且每个幂零
       矩阵都与一个标准形的矩阵相似。

       看出这一点的另一种方式是注意：幂零矩阵在若干次迭代后把所有
       向量都送到零，但这与在特征子空间上的作用 $\vec{v}\mapsto \lambda\vec{v}$
       相矛盾，除非
       $\lambda$ 为零。
     \end{answer}
   \item
     在二维空间上存在指数为三的幂零变换吗？
     \begin{answer}
       不存在，由 \nearbylemma{le:RangeAndNullChains}，对二维
       空间上的映射，零度
       在第二次迭代时就会增长到尽可能大。
     \end{answer}
   \item
     在 \nearbytheorem{th:NilMapHasStrBas} 的证明中，为什么归纳基础的情形不是幂零指数为零？
     \begin{answer}
       一个变换的幂零指数为零，仅当开启链的
       向量
       必须是 $\zero$ 时，也就是说，仅当 $V$ 是平凡空间时。
     \end{answer}
   \recommended \item
     设 \( \map{t}{V}{V} \) 是线性变换，并设
     \( \vec{v}\in V \) 满足 \( t^k(\vec{v})=\zero \) 但
     \( t^{k-1}(\vec{v})\neq\zero \)。
     考虑 $t$-链
     $\sequence{\vec{v},t(\vec{v}),\dots,t^{k-1}(\vec{v})}$。
     \begin{exparts}
       \partsitem 证明 \( t \) 是该链中向量组张成的空间上的变换，
         也就是说，证明 \( t\) 限制在该张成空间上的值域
         是该张成空间的子集。
         我们说这个张成空间是 \definend{\( t \)-不变的}
         子空间。\index{invariant subspace}
       \partsitem 证明该限制是幂零的。
       \partsitem 证明该 $t$-链
         线性无关，从而是其张成空间的基。
       \partsitem 用该 $t$-链基表示这个限制映射。
     \end{exparts}
     \begin{answer}
       \begin{exparts}
         \partsitem 张成空间中的任意成员 $\vec{w}$ 都是
           一个线性组合
           $\vec{w}=c_0\cdot \vec{v}+c_1\cdot t(\vec{v})+\dots
            +c_{k-1}\cdot t^{k-1}(\vec{v})$。
           于是，由映射的线性性，
           $t(\vec{w})=c_0\cdot t(\vec{v})+c_1\cdot t^2(\vec{v})+\dots
            +c_{k-2}\cdot t^{k-1}(\vec{v})+c_{k-1}\cdot \zero$
           也在这个张成空间中。
         \partsitem 前一小题中的操作迭代 $k$ 次后，
           会得到零的线性组合。
         \partsitem 若 \( \vec{v}=\zero \)，则该集合是空集，按定义
           线性无关。
           否则，写 \( c_1\vec{v}+\dots+c_{k-1}t^{k-1}(\vec{v})=\zero \)
           并对两边施加 \( t^{k-1} \)。
           右边得到 \( \zero \)，而左边得到
           \( c_1t^{k-1}(\vec{v}) \)；于是 \( c_1=0 \)。
           如此继续，对两边施加 \( t^{k-2} \)，
           依此类推。
         \partsitem 当然，$t$ 把这条长为 $k$ 的单一 $t$-链作为基作用在其张成空间上。
           \begin{equation*}
             \begin{mat}
               0 &0 &0 &0 &\ldots &0 &0 \\
               1 &0 &0 &0 &\ldots &0 &0 \\
               0 &1 &0 &0 &\ldots &0 &0 \\
               0 &0 &1 &0 &       &0 &0 \\
                 &  &  &\ddots          \\
               0 &0 &0 &0 &       &1 &0 \\
             \end{mat}
           \end{equation*}
       \end{exparts}
     \end{answer}
   \item \label{exer:NilMapWillHaveStrBas}
     完成 \nearbytheorem{th:NilMapHasStrBas} 的证明。
     \begin{answer}
       我们必须证明
       \( B\union\hat{C}\union\set{\vec{v}_1,\dots,\vec{v}_j} \) 线性
       无关，其中 \( B \) 是
       \( \rangespace{t} \) 的 \( t \)-链基，\( \hat{C} \) 是
       \( \nullspace{t} \) 的基，且
       \( t(\vec{v}_1)=\vec{\beta}_1,\dots,t(\vec{v}_i)=\vec{\beta}_i \)。
       写出
       \begin{equation*}
         \zero=c_{1,-1}\vec{v}_1+c_{1,0}\vec{\beta}_1
                +c_{1,1}t(\vec{\beta}_1)+\dots+
                c_{1,{h_1}}t^{h_1}(\vec{\vec{\beta}}_1)
                +c_{2,-1}\vec{v}_2+\dots+c_{j,h_i}t^{h_i}(\vec{\beta_i})
       \end{equation*}
       并施加 \( t \)。
       \begin{multline*}
         \zero=c_{1,-1}\vec{\beta}_1+c_{1,0}t(\vec{\beta}_1)
                +\dots+
                c_{1,h_1-1}t^{h_1}(\vec{\vec{\beta}}_1)+c_{1,h_1}\zero      \\
             +c_{2,-1}\vec{\beta}_2+\cdots+c_{i,h_i-1}t^{h_i}(\vec{\beta_i})
             +c_{i,h_i}\zero
       \end{multline*}
       由于 \( B\union\hat{C} \)
       是一组基，可断定系数 \( c_{1,-1},\dots,c_{1,h_i-1},
       c_{2,-1},\dots,c_{i,h_i-1} \) 全为零。
       代回第一个行间公式，可知
       其余系数也为零。
     \end{answer}
   \item \label{exer:MatNilIffMapNil}
     证明 \nearbydefinition{def:nilpotent} 中给出的「幂零变换」与「幂零矩阵」这两个说法彼此吻合：~
     一个映射是幂零的当且仅当它由一个
     幂零矩阵表示。
     （是说一个变换是幂零的当且仅当存在一组基，使得该映射关于该基的表示是
     幂零矩阵，还是说任何表示都是幂零矩阵？）
     \begin{answer}
       对任意基 $B$，
       变换~$n$ 是幂零的当且仅当
       $N=\rep{n}{B,B}$ 是幂零矩阵。
       这是因为只有零矩阵才表示零映射，
       于是 \( n^j \) 是零映射当且仅当 \( N^j \)
       是零矩阵。
     \end{answer}
   \item
     设 \( T \) 是指数为四的幂零变换。
     \( T^3 \) 的值域空间最大能有多大？
     \begin{answer}
       它可以是任何大于或等于一的维数。
       要构造一个指数为四的幂零变换，
       其立方有维数为~$k$ 的值域空间，取一个向量空间、
       该空间的一组基，以及一个按如下方式
       作用于这组基的变换。
       \begin{equation*}
          \begin{strings}{ccccccccc}
             \vec{\beta}_1 &\mapsto &\vec{\beta}_2 &\mapsto &\vec{\beta}_3
               &\mapsto &\vec{\beta}_4 &\mapsto &\zero  \\
             \vec{\beta}_5 &\mapsto &\vec{\beta}_6 &\mapsto &\vec{\beta}_7
               &\mapsto &\vec{\beta}_8 &\mapsto &\zero  \\
             &\vdotswithin{\vec{\beta}_{4k-3}}                       \\
             \vec{\beta}_{4k-3} &\mapsto &\vec{\beta}_{4k-2}
               &\mapsto &\vec{\beta}_{4k-1}
               &\mapsto &\vec{\beta}_{4k} &\mapsto &\zero  \\
             &\vdotswithin{\vec{\beta}_{4k-3}}              \\
             \multicolumn{8}{l}{%
               \text{\textit{--可能还有别的更短的链--}}}
          \end{strings}
       \end{equation*}
       所以 $T^3$ 的值域空间的维数可以任意大。
       它最小为一\Dash 这里
       必须至少有一条链，否则该映射的幂零指数
       就不会是四。
     \end{answer}
   \item
     回忆相似矩阵有相同的特征值。
     证明其逆命题不成立。
     \begin{answer}
       下面两个矩阵的特征值都只有零
       \begin{equation*}
         \begin{mat}
           0  &0  \\
           0  &0
         \end{mat}
         \qquad
         \begin{mat}
           0  &0  \\
           1  &0
         \end{mat}
       \end{equation*}
       但它们并不相似（它们有不同的标准形
       代表元，即它们自身）。
     \end{answer}
   \item \nearbylemma{lem:RestONeToOne} 表明，对任意线性
     变换 \( \map{t}{V}{V} \)，
     限制 \( \map{t}{\genrangespace{t}}{\genrangespace{t}} \)
     都是单射。
     证明它也是满射，从而是一个自同构。
     它必须是恒等映射吗？
     \begin{answer}
       由 \nearbylemma{le:RangeAndNullChains}，它是满射。
       它不必是恒等映射：考虑这个映射 $\map{t}{\Re^2}{\Re^2}$。
       \begin{equation*}
         \colvec{x \\ y}\mapsunder{t}\colvec{y \\ x}
       \end{equation*}
       对该映射有 $\genrangespace{t}=\Re^2$，而 $t$ 不是恒等映射。
     \end{answer}
   \item
     证明幂零矩阵相似于一个除上对角线上的 $1$ 组成的块以外全为零的矩阵。
     \begin{answer}
       对链基作一个简单的重排即可。
       例如，与这个链基相关的映射
       \begin{equation*}
          \begin{strings}{ccccccc}
             \vec{\beta}_1 &\mapsto &\vec{\beta}_2 &\mapsto &\zero
          \end{strings}
       \end{equation*}
       关于
       $B=\sequence{\vec{\beta}_1,\vec{\beta}_2}$ 由这个矩阵表示
       \begin{equation*}
         \begin{mat}
           0  &0  \\
           1  &0
         \end{mat}
       \end{equation*}
       但关于
       $B=\sequence{\vec{\beta}_2,\vec{\beta}_1}$ 却这样表示。
       \begin{equation*}
         \begin{mat}
           0  &1  \\
           0  &0
         \end{mat}
       \end{equation*}
     \end{answer}
   \recommended \item
     证明若一个变换的值域空间与零空间相同，
     则其定义域的维数是偶数。
     \begin{answer}
       设 $\map{t}{V}{V}$ 是该变换。
       若 \( \rank (t)=\nullity (t) \)，则等式
       \( \rank(t)+\nullity(t)=\dim (V) \) 表明
       \( \dim (V) \) 是偶数。
     \end{answer}
   \item
     证明若两个幂零矩阵可交换，则它们的乘积与
     和也是幂零的。
     \begin{answer}
       矩阵要是幂零的，它们必须是方阵。
       要可交换，它们必须大小相同。
       于是它们的乘积与和都有定义。

       把这两个矩阵称为 \( A \) 与 \( B \)。
       要看出 \( AB \) 是幂零的，计算
       $
          (AB)^2=ABAB=AABB=A^2B^2$，
       以及
          $(AB)^3=A^3B^3$，依此类推，
       又因 \( A \) 是幂零的，该乘积最终为零。

       和的情形类似；利用二项式定理。
      \end{answer}
   % 2014-Dec-19 These two questions and answers are fine; I commented them to make the pages fit better.
   % \item
   %   Consider the transformation of \( \matspace_{\nbyn{n}} \) given
   %   by \( t_S(T)=ST-TS \) where \( S \) is an \( \nbyn{n} \) matrix.
   %   Prove that if \( S \) is nilpotent then so is \( t_S \).
   %   \begin{answer}
   %     Some experimentation gives the idea for the proof.
   %     Expansion of the second power
   %     \begin{equation*}
   %       t^2_S(T)=S(ST-TS)-(ST-TS)S=S^2-2STS+TS^2
   %     \end{equation*}
   %     the third power
   %     \begin{align*}
   %       t^3_S(T)
   %         &=S(S^2-2STS+TS^2)-(S^2-2STS+TS^2)S  \\
   %         &=S^3T-3S^2TS+3STS^2-TS^3
   %     \end{align*}
   %     and the fourth power
   %     \begin{align*}
   %       t^4_S(T)
   %         &=S(S^3T-3S^2TS+3STS^2-TS^3)-(S^3T-3S^2TS+3STS^2-TS^3)S  \\
   %         &=S^4T-4S^3TS+6S^2TS^2-4STS^3+TS^4
   %     \end{align*}
   %     suggest that the expansions follow the Binomial Theorem.
   %     Verifying this by induction on the power of $t_S$ is routine.
   %     This answers the question because, where the index of nilpotency of
   %     $S$ is $k$, in the expansion of $t^{2k}_S$
   %     \begin{equation*}
   %       t^{2k}_S(T)=\sum_{0\leq i\leq 2k}(-1)^i\binom{2k}{i} S^iTS^{2k-i}
   %     \end{equation*}
   %     for any $i$ at least one of the $S^i$ and $S^{2k-i}$
   %     has a power higher than $k$, and so the term gives the zero matrix.
   %   \end{answer}
   % \item
   %   Show that if \( N \) is nilpotent then \( I-N \) is
   %   invertible.
   %   Is that `only if' also?
   %   \begin{answer}
   %     Use the geometric series:
   %     $
   %        I-N^{k+1}=(I-N)(N^k+N^{k-1}+\cdots+I)
   %     $.
   %     If \( N^{k+1} \) is the zero matrix then we have a right inverse for
   %     \( I-N \).
   %     It is also a left inverse.

   %     This statement is not `only if' since
   %     \begin{equation*}
   %        \begin{mat}[r]
   %           1  &0  \\
   %           0  &1
   %        \end{mat}
   %        -\begin{mat}[r]
   %           -1  &0  \\
   %           0  &-1
   %        \end{mat}
   %     \end{equation*}
   %     is invertible.
   %    \end{answer}
 %  \recommended \item
 %    Show that when a nilpotent matrix \( T \) is in canonical form
 %    then the number of
 %    blocks of \( \nbyn{(k+1)} \) matrices that are all zeros
 %    except for subdiagonal ones is
 %    \( 2\,(\nullity (T^k))
 %      -\nullity (T^{k+1})-\nullity (T^{k-1}) \).
\index{nilpotent|)}
\end{exercises}
